\documentclass[12pt]{book}
\usepackage{amsmath}
\usepackage{amssymb}
\usepackage{geometry}
\geometry{margin=1in}

\title{Volume 11: Validation, Benchmarks, and Empirical Tests \\ Book 02: SDT vs Standard Physics}
\author{James C. Harvey}
\date{December 2025}

\begin{document}
\maketitle
\tableofcontents

\chapter{Comparative Framework}

\section*{Abstract}
This book compares SDT predictions with standard physics results, emphasizing where SDT derives
quantities geometrically rather than by postulate.

\section{Comparison Principle}
The comparison is performed at the level of derived constants and functional forms. SDT does not
introduce independent fields where geometry suffices.

\chapter{Functional Form Comparison}

\section*{Abstract}
Comparisons focus on whether SDT reproduces the same functional relationships without external postulates.

\section{Form Mapping}
If a standard formula is $F = f(r)$, SDT seeks $F = f(\kappa(r))$ with $\kappa$ derived from boundary
geometry and occlusion.

\chapter{Parameter Discipline}

\section*{Abstract}
SDT avoids adjustable parameters. This chapter formalizes the rule.

\section{No Free Fits}
If a constant is not derivable from $P_{\text{CMB}}$, $K_{\text{bulk}}$, $c$, or geometry, it is not
introduced.
\end{document}
